% !PRINTED 72
\weberpage{72}
%<<CONT
and thus by 3. a divisor $\Omega$ of $\Omega_m$ will fail to be primary
if and only if one of the groups $Q_1, Q_2, \ldots Q_\tau$ is contained
in the group $\mathfrak{A}$ belonging to $\Omega$. If, then, we call
those divisors of the group $\mathfrak{N}$ which contain none of the groups $Q_1, Q_2 \ldots Q_\tau$
as divisors primary divisors, we can state the theorem:

\gesperrt{5. To obtain all primary divisors of $\Omega_m$,
one has to seek out all primary divisors of the group $\mathfrak{N}$
and to form the associated fields.}

\gesperrt{6. But if one sets up all primary divisors of all
full cyclotomic fields, one obtains
every cyclotomic field, and each one
only once, and indeed each represented by
roots of unity of the lowest possible degree.}

The field $\Omega_{2m}$, when $m$ is odd, has no primary divisor at all and is identical with $\Omega_m$. In all other cases
$\Omega_m$ is at least its own primary divisor.

For what follows it is further necessary that we become clear
about the extent and the constitution of the groups $Q$.

\begin{center}
\S.\,19.

\gesperrt{Primary and non-primary divisors of the group $\mathfrak{N}$.}
\end{center}

We denote by $q$ any one of the primes dividing $m$
and set $m = q m'$. We then consider the
group
\[
Q \equiv 1\,(\mathrm{mod}\ m').
\]

To determine the conditions for a number in $Q$, we apply
the representation of the group $\mathfrak{N}$ by a basis and the
notation of the indices, as we introduced and explained them
in §. 16, and accordingly denote the group reciprocal to $Q$
by $P$.

The indices of a number in $Q$ we denote by
\[
\gamma_{-1}, \gamma_0, \gamma_1, \ldots \gamma_\mu,
\]
and those of a number in $P$ by
\[
\delta_{-1}, \delta_0, \delta_1, \ldots \delta_\mu.
\]
